\section{Finite cocycles and integral crystalline period representatives}

% BEGIN P85C1
\paragraph{P85C1. The two integral period lattices.}
Use the field, character set and test lattices of Q84T1--2. Write $p$ for the
prime, $[-b,H]$ for the height interval, $e=p^a-1$ and $E=u^e+p$.
Use R85P1--3 to supply a realized module. This has the actual $\mathrm{Pack}$ and
$\mathrm{GalDat}$ of the pin, including their Witt identification,
comparison on the common integral lattice, tame descent and horizontal
periods. Write $M$ for the integral lattice, $L$ for the span of its
Frobenius columns, and $\mathcal L$ for the extension of $L$ to Breuil's
divided-power ring $S$. Thus $E^bL\subset M$ and $E^HM\subset L$.
Let $W=W(C_p^\flat)$ and $A_{\inf}=W(\mathcal O_{C_p^\flat})$;
extend both by the coefficient integer ring $O$. Put $\mu=[\epsilon]-1$,
$\nu=[\epsilon^{1/p}]-1$. Evaluation at the compatible uniformizers gives
$\varphi\nu=\mu$ and $E([\pi^\flat])=v\mu/\nu$ with $v$ a unit
of $A_{\inf}$, since both sides generate $\ker\theta$.
Here $\theta\nu=\zeta_p-1\ne0$, so $\nu$ is a unit in $B_{\rm dR}^+$.
Both $\mu$ and $\nu$ are units in $B_{\rm cris}$ by the cyclotomic
period identities.

In the Witt coordinates let $A$ be the Frobenius matrix of $M$.
Both $E^bA$ and $E^H A^{-1}$ are integral. An equation
$\varphi z=Cz+d$ with $C,d$ over $A_{\inf}$ forces every
Witt solution $z\in W^n$ to be over $A_{\inf}$. Modulo $p$ a
negative minimum coordinate valuation $s$ would give $ps<s$
on the left and valuations $\ge s$ on the right. Lift an integral
residue, subtract it and divide the equation by $p$. The new
inhomogeneous term is integral since $pW\cap A_{\inf}=pA_{\inf}$.
Induct on the Witt digits. This proof allows $O$-coefficients by
choosing an integral $\mathbb Z_p$-basis of $O$, with Frobenius
acting trivially on those coefficients.

In a basis of $T$ the Witt realizations give columns of a
matrix $U$, with $A\varphi U=U$. The inverse matrix $V$ gives
the comparison to that basis, and $\varphi V=VA$. The preceding
test on $\nu^H U$ and the rows of $\nu^b V$ gives integrality:
their Frobenius equations have coefficients unit multiples of
$E^H A^{-1}$ and $E^b A$, respectively. Taking Frobenius,
$\mu^HU$ belongs to the evaluated $L$ and $\mu^bVA$ is integral.
Thus $\nu^b M_{A_{\inf}}$ and $\mu^b L_{A_{\inf}}$ map into
$A_{\inf}\otimes T$, while $\nu^H T$ and $\mu^HT$ map into the
corresponding lattices. Extend these comparisons to $A_{\rm cris}$.

For $E^b y\in M$ solve $A\varphi m-m=y$ over $W\otimes O$.
In the $T$ basis this is $\varphi-1$, surjective by
Artin--Schreier surjectivity on $C_p^\flat$ and successive
$p$-adic lifting, with fixed lattice $T$. Set $z=\nu^{H+b}m$.
In $\varphi z=\mu^{H+b}A^{-1}(m+y)$ the coefficient of $z$
and the inhomogeneous term are integral by $E^HA^{-1}$ and
$E^b y$. Thus $z$ is integral. Moreover $\mu^{H+b}\Phi m$
belongs to the integral Frobenius-column lattice, where
$\Phi m=A\varphi m$. Consequently $\Phi m$ has a crystalline
value $b_0$, obtained by evaluating common integral coordinates
and dividing by $\mu^{H+b}$, and
$\mu^{H+2b}b_0\in A_{\inf}\otimes T$. The Witt vector $m$ itself
has a crystalline value $b_1$ with
$\nu^{H+2b} b_1\in A_{\inf}\otimes T$ by the $M$ estimate.
For the definition using the Frobenius columns, multiply
these latter common coordinates further by $\mu^H$ to enter
$L_{A_{\inf}}$ on both sides. This preserves the value.
In particular $b_1\in B_{\rm dR}^+\otimes T$. Uniqueness of
crystalline values follows by clearing both denominators into
$A_{\inf}$, using its injection into $W$, and evaluating in
$B_{\rm cris}$.
% END P85C1

% BEGIN P85C2
\paragraph{P85C2. The finite class supplies the corrected invariant vector.}
Let $c\in Z_f^1(T)$ for a good lattice with realized tame model.
Take $T^+=T\oplus O v$ with $gv=v+c(g)$, and its quotient $O$.
A crystalline primitive for $c$ makes $v$ minus that primitive
invariant over $B_{\rm cris}$. Together with a crystalline basis
for $T$ over $K_t$, it gives all periods for the extension there.
The new line is trivial, hence its inertia character and its
filtration jump add $0$ to lists already containing $0$, by
strictness on the de Rham extension and semisimplicity of finite
period inertia in characteristic zero. Thus the same good-model
construction applies to the extension and unit. Denote the two
realized maps by $F,G$. The unit has its invariant generator $u_0$
with $\Phi u_0=u_0$ evaluating to $1$: the explicit unit model
maps to the chosen model by integral full faithfulness, and
descent invariance follows by evaluating on Witt coordinates.

Let $w=2(N-1)R$. Use the sharp cokernel bound on invariant lattices
of E85C4. Take
$x\in(M^+)^\Delta$ with $Gx=p^w u_0$. Then $G(\Phi x-x)=0$,
and $E^b(\Phi x-x)\in M^+$. By Q84T3 there is a unique
$y\in E^{-b}M$ such that $Fy=\Phi x-x$. Uniqueness gives
$\Delta$-invariance. The map on analytic principal parts is
injective: after inverting $p$ the model sequence splits as a
sequence of underlying modules since the terminal module is
free and its finite-length cokernel is killed. The resulting
left inverse is analytic on the $E$-disc and detects regularity.
Since $\mathrm{Mon}(\Phi x-x)=-[E\nabla x]=0$ by preservation of
the physical analytic lattice under $E\nabla$, we have
$y\in Y$ and $\mathrm{Mon}(y)=0$.

Write $\Xi(d)$ for the Frobenius-compatible horizontal section
with constant term $d$. Fix a disc of radius strictly between
$p^{-1/(e(p-1))}$ and $p^{-1/(ep)}$. Apart from the first
matrix $A(u)$ all factors $A(u^{p^j})$, $j\ge1$, and their
inverses are analytic there. The normalized Frobenius product
given by iterating
$\Xi_\tau=A_\tau\varphi(\Xi_{\tau-1})A_\tau(0)^{-1}$ from $1$
converges there after removing the first factor: the error from
a new tail factor is $O(r^{p^j})$ and conjugation by the
constant matrices and their inverses costs at most exponential
growth in $j$. The inverse product likewise makes
$A^{-1}\Xi$ invertible analytic there. Its coefficients and those
of its inverse belong to $S[1/p]$: in the $u^n$ coefficient test
the weight is $|\lfloor n/e\rfloor!|_p$, of $n$th-root limit
$p^{-1/(e(p-1))}$. Thus the horizontal and Frobenius-column
frames agree over $S[1/p]$. Near a root of $E$ the physical and
horizontal frame change and its inverse have finite poles, bounded
by the two height inequalities. Functoriality follows by applying
the Frobenius recurrence to a map and using equality at zero.

Put $d=\mathrm{red}(\Phi x)$. The image of $\Phi x-\Xi(d)$
under $G$ is zero, since $\Xi(\mathrm{red}\,u_0)=u_0$ by the
same recurrence. In Frobenius-column bases after inverting $p$,
the sequence stays exact: its matrices are Frobenius pullbacks
of the split sequence on $M$. Thus there is a unique
$\sigma\in\mathcal L[1/p]$ with $F\sigma=\Phi x-\Xi(d)$.
It is invariant. The identity $\zeta(\Phi x-x)=\nabla(\Phi x)$
gives $\nabla\sigma=\zeta y$.

Choose a Witt solution $m$ as in P85C1. The defect
$Fm-(x-p^wv)$ is Frobenius fixed and killed by $G$, hence
is a member of the actual kernel lattice $T$. Subtract its
preimage to arrange $Fm=x-p^wv$. Compare crystalline
values of the Frobenius of this equality. Writing $j$ for
crystalline evaluation in the Frobenius-column frame gives
\[
   F\big(j(\sigma)-b_0\big)=p^w(1\otimes v)-j(\Xi(d)).
\]
The horizontal period on the right is fixed by $G_{K_t}$
from the realized model and by the tame covers from
$\Delta$-invariance. These groups generate the absolute
Galois group as detailed below. Take Galois differences and use
the split inclusion after tensoring with $B_{\rm cris}$. This gives
$(g-1)(j(\sigma)-b_0)=1\otimes p^w c(g)$. Together with the
Witt equation and $\nabla\sigma=\zeta y$ this is the native
relation $\mathrm{ClassOf}(y,p^w c)$.
% END P85C2

% BEGIN P85C3
\paragraph{P85C3. The rational crystalline primitive and the de Rham test.}
For any $y\in Y$, form the extra Frobenius column $q+y$ on
$M\oplus O[[u]]q$, with $q$ fixed by tame descent. This auxiliary
matrix is invertible after inverting $E$. The normalized horizontal
section of P85C2 applies also to this upper triangular matrix.
The difference $\Phi q-\Xi^+(\mathrm{red}\,\Phi q)$ has last
coordinate zero, so defines $\sigma\in\mathcal L[1/p]$.
Indeed both terms belong to the extended Frobenius-column lattice
and its intersection with last coordinate zero is the old lattice
by the block matrix. The vector is descent invariant. The canonical
connection is characterized in the horizontal frame by
$N=u\,d/du$. Its identity $\nabla\Phi=p\Phi\nabla$ gives
$\nabla\Phi q=p\Phi(\nabla\Phi q-\nabla y)$. Its zero-constant
solution is $\zeta y=-\sum_{n\ge1}(p\Phi)^n\nabla y$:
the difference of two solutions has zero constant and is successively
divisible by arbitrarily high powers of $u$. Thus
$\nabla\sigma=\zeta y$. Choose $m,b_0,b_1$ as in P85C1. Then
$B=j(\sigma)-b_0$ is the pin's potential and $b_y=b_0-b_1$ is the
crystalline value of $y$. Taking a common denominator shows that
evaluation of a Witt vector with a crystalline value commutes with
$\Phi$; this follows on the common integral lattice from
$A\varphi U=U$ in P85C1 and after applying Frobenius to the
denominator. In particular $\varphi b_0-b_0=j(\Phi y)$.
Here $\Phi y\in\mathcal L[1/p]$ since $\varphi E/p$ is a unit in
$S$. The identity $\nabla(\Phi\sigma-\sigma-\Phi y)=0$
follows from the recurrence for $\zeta$. Hence
$(\varphi-1)B$ is the evaluation of a constant horizontal
section. It is fixed by $G_{K_t}$. Consequently $(h-1)B$ is
Frobenius fixed for $h\in G_{K_t}$.

Suppose now $\mathrm{Mon}(y)=0$. In the rational germ at
$u=\pi_0$ write $z=\log(u/\pi_0)$ and work separately in the
unramified coefficient components. Extend constants to $K_t$ as
needed. The Euler operator is $\partial_z$ in the horizontal
frame and $E/z$ is a unit. Write $\widehat M$ for the physical
lattice over the formal $z$-disc. It lies between
$z^d$ and $z^{-d}$ times the horizontal power-series lattice
for some $d$, by P85C2, and is stable under $z\partial_z$
because $E\nabla$ preserves the physical analytic lattice.
The vector $s=\sigma-y$ has finite poles and
$E\nabla s=E\zeta y-E\nabla y$ is in $\widehat M$ by
$\mathrm{Mon}(y)=0$. Increase $d$ so that $s$ lies in the
larger lattice. Modulo $z^d$ times the horizontal lattice the
operator $z\partial_z$ is diagonal with distinct integral
eigenvalues from $-d$ to $d-1$. Its stable subspace from
$\widehat M$ is a sum of intersections with these eigenspaces,
as is seen by applying the polynomial spectral projectors.
It follows from $z\partial_z s\in\widehat M$ that all nonconstant
terms of $s$ modulo this truncation are in $\widehat M$.
Thus $s$ differs from a constant horizontal vector by a member
of $\widehat M$. Only coefficients in $K_t$ with the original
external coefficient scalars are involved: all Taylor coefficients
are obtained from convergent power series at the algebraic point
$\pi_0$ in this complete field.

Under evaluation in $B_{\rm dR}$ the germ substitution is
$u=[\pi^\flat]$, where $[\pi^\flat]-\pi_0\in\ker\theta$.
A regular physical germ evaluates in $B_{\rm dR}^+\otimes T$.
Indeed the physical basis takes values there by the
$\nu^bM$ estimate in P85C1, and a formal series in the last
difference converges in the filtration topology. The assertion
continues to hold for each indicated coefficient embedding and
agrees with $j$ on $\mathcal L[1/p]$ and the crystalline value on
$E^{-b}M$: clear $E$ denominators, use evaluation of common integral
coordinates, and pass to the Taylor expansion. For the latter
passage polynomials are dense in the weighted $S$ norm. On any
fixed Taylor quotient evaluation is continuous, since the radius
$p^{-1/e}$ at $\pi_0$ is strictly smaller than
$p^{-1/(e(p-1))}$. The evaluated horizontal constant is fixed
by $G_{K_t}$, including its $K_t$ scalars. It follows that
$(h-1)(j(\sigma)-b_y)$ has filtration at least zero. The same
holds for $b_1$ by P85C1. Since $B=j(\sigma)-b_y-b_1$,
Fontaine's identity $B_{\rm cris}^{\varphi=1}\cap
B_{\rm dR}^+=\mathbb Q_p$ gives $(h-1)B\in1\otimes V$.
Apply it on a $\mathbb Q_p$-basis, so ramification of the
coefficient field introduces no restriction.
% END P85C3

% BEGIN P85C4
\paragraph{P85C4. The exact integral return exponent.}
Suppose $\mathrm{Mon}(y)=0$, $p^c\zeta y\in\mathcal L$ and
$8H+1<p$, with $1\le b<H$. First take $h\in G_{K_t}$.
The ratio of $h[\pi^\flat]$ to $[\pi^\flat]$ is
$1+D=[\epsilon]^k$ for some $k\in\mathbb Z_p$, so $D\in
\mu A_{\inf}$ (equivalently $[\epsilon]^k-1$ and $\mu$ differ
by a factor in $A_{\inf}$, by the cyclotomic tower).
In the horizontal frame which $h$ fixes, the Taylor formula is
\[
 p^c(h-1)j(\sigma)=\sum_{n\ge1}
       \frac{D^n}{n!}\,j\big(p^c\nabla(\nabla-1)\cdots
                       (\nabla-n+1)\sigma\big).
\]
Every vector on the right before evaluation lies in $\mathcal L$,
since its first derivative is integral after multiplication by $p^c$
and $\nabla$ preserves that lattice. The coefficients $D^n/n!$
are integral in $A_{\rm cris}$ and tend to zero $p$-adically:
write $\mu=\nu\xi$, $\xi$ a divided-power generator, and observe
$\nu^{p-1}\equiv\xi\pmod p$ and $\xi^p\in pA_{\rm cris}$.
For a scalar polynomial the formula is the binomial expansion
at $(1+D)u$. For coefficients in $S[1/p]$ it follows by weighted
polynomial approximation: clear a common $p$-denominator in
the horizontal scalar coordinates and in the finitely many
horizontal basis columns. The falling Euler derivatives of
$u^j$ have integral coefficients relative to its weighted norm,
uniformly for each derivative. The divided-power coefficients
above make the tails converge uniformly. This proves the identity
in the crystalline lattice after the common denominator is cleared,
and the integral right side already has Frobenius-column
coordinates over $A_{\rm cris}$. P85C1 now gives
$\mu^b p^c(h-1)j(\sigma)\in A_{\rm cris}\otimes T$.
Also $\mu^{H+2b}(h-1)b_0\in A_{\inf}\otimes T$:
the ideal $(\mu)$ is Galois stable since $h\epsilon=\epsilon^{\chi_p(h)}$
with $\chi_p(h)\in\mathbb Z_p^\times$. It follows that
$\mu^{H+2b}p^c(h-1)B\in A_{\rm cris}\otimes T$.

The loss from $\mu$ can be removed on the rational subspace in P85C3.
In fact for $0\le j<p-1$, $\mu^j\notin p A_{\rm cris}$.
Map $A_{\inf}$ to $\mathcal O_{C_p^\flat}/(\bar\xi^p)$ by reducing
modulo $p$, where the bar denotes the zeroth Witt coordinate.
On the ideal $(\bar\xi)$ in this quotient use the divided powers
$x^n/n!$ for $n<p$, zero for higher indices. The divided-power
identities hold since products of total degree $\ge p$ vanish.
This extends to the divided-power envelope and its $p$-adic
completion. In the tilt valuation normalized by
$v(\bar\xi)=1$ one has $v(\bar\mu)=p/(p-1)$: the valuations
of $\zeta_{p^n}-1$ are $1/(p^{n-1}(p-1))$, and taking the tilt
multiplies them by $p^n$. Thus $v(\bar\mu^j)<p$ for $j<p-1$,
so its image is nonzero. For a rational lattice coordinate $t$
with $\mu^j t\in A_{\rm cris}$, negative $p$-valuation of $t$
would imply $\mu^j\in pA_{\rm cris}$, an impossibility.
Use a $\mathbb Z_p$-basis of $T$ and $j=H+2b<p-1$.
We obtain $p^c(h-1)B\in1\otimes T$.

It remains to cover the whole Galois group. Put
$K_\infty=\bigcup_n\mathbb Q_p(\pi_n^e)$ and
$K_{t,\infty}=\bigcup_n\mathbb Q_{p^a}(\pi_n)$. Each
automorphism of $K_t/\mathbb Q_p$ extends compatibly to
$K_{t,\infty}$ fixing each $\pi_n^e$: send $\pi_n$ to
$\zeta_n\pi_n$ with $\zeta_n\in\mu_e$ and
$\zeta_{n+1}^p=\zeta_n$. The unramified part acts on constants.
The relations suffice since the field tower has defining
Eisenstein polynomials $X^{e p^n}+p$ at each finite level.
Extend to the algebraic closure. This gives $g=hk$ with
$h\in G_{K_t}$ and $k$ preserving the tower and inducing such
a tame action. Such $k$ fixes the evaluations of the invariant
vectors $y,\sigma$ by the realized descent comparison, also on
$S[1/p]$ by scalar extension and continuity. The Witt difference
$km-m$ is therefore Frobenius fixed, hence an actual vector
$t\in T$. Uniqueness of crystalline values gives
$kb_0-b_0=1\otimes t$. Thus $(k-1)B=-1\otimes t$ is already
integral. The identity $(hk-1)B=h(k-1)B+(h-1)B$ proves the
same $p^c$ bound on every $g$.

Write $1\otimes\kappa(g)=p^c(g-1)B$. The cocycle law is immediate.
Continuity follows from the continuous action on $B_{\rm dR}$ and
the induced usual topology on the finite-dimensional rational
subspace: embed in a fixed $t^{-d}B_{\rm dR}^+$ step, multiply
by $t^d$, then read the degree-$d$ coefficient on the rational
subspace in its filtration quotient. The primitive $p^c B$
belongs to $B_{\rm cris}\otimes V$, hence $\kappa$ is finite.
Scaling the Witt and crystalline data by $p^c$ supplies the native
$\mathrm{ClassOf}(p^cy,\kappa)$.
% END P85C4

% BEGIN P85C5
\paragraph{P85C5. The return exponent for a monodromy-zero pair.}
Put $c_1=b(R+1)$ and let $y=y'+\Phi m$ with $y'\in Y'$,
$m\in M^\Delta$, $\mathrm{Mon}(y)=0$. Then $p^{c_1}\zeta y$
belongs to the full lattice $\mathcal L$. We give the argument
because an estimate on a class modulo analytic vectors alone
would not establish the hypothesis of P85C4.

Put $q=E\zeta y-E\nabla y$, an element of $M_{\rm an}$ by zero
monodromy. We claim there are $P\in L$ and
$F\in\mathrm{Fil}^{10H}\mathcal L$ with
\[
        p^{c_1}E^{2H}q=P+p^{-3H}F.
\]
For the first part of $E^{2H+1}\zeta y$, use M85N4 on $\zeta y'$:
its error is in $p^{-3H}\mathrm{Fil}^{10H}\mathcal L$ and its
other summand in $L$, since $p>14H$. Multiplication by $p^{c_1}$
and $E^{2H+1}$ preserves both bounds. For the second, M85N5
gives $p^{c_1}\zeta(\Phi m)\in\mathcal L$; split it into a term
in $L$ and a term in $\mathrm{Fil}^{10H}\mathcal L$ by M85N1.
Finally $E^by\in M$ by the upper height bound on $\Phi m$,
and the Leibniz calculation in M85N1 splits $E^{2H+1}\nabla y$
with the same bounds. Subtracting proves the claim.

The term $E^{-2H}F$ belongs to $M_{\rm an}$ by M85N1. Hence
$X=E^{-2H}P$ belongs to $M_{\rm an}$ as well, using the displayed
identity and regularity of $q$. Also $E^{2H+b}X\in M$ since
$E^bL\subset M$. Divide that numerator by the distinguished
polynomial $E^{2H+b}$ in each coordinate. Analytic regularity at
all roots of $E$ forces the remainder to vanish, so $X\in M$.

Write $a_0=p/\varphi(E)$, a unit in $S$ by M85N1.
For $j=10H<p$, Frobenius sends $\mathrm{Fil}^jS$ into
$p^jS$: on $E^i/i!$, $i\ge j$, the valuation of the extracted
factor is $i-v_p(i!)\ge j$. This last bound is immediate for
$j\le i<p$; for $i\ge p$ it follows from
$v_p(i!)\le(i-1)/(p-1)$ and integrality, giving $i-v_p(i!)\ge p-1$.
Continuity gives the same on the closed ideal. Since
$p^b\Phi\mathcal L\subset\mathcal L$, write
$\Phi F=p^{10H-b}F'$ with $F'\in\mathcal L$.
The recurrence for $\zeta$ gives $\zeta y=a_0\Phi q$. Applying
Frobenius to the displayed decomposition yields
\[
 p^{c_1}\zeta y
     =a_0\big(\Phi X+a_0^{2H}p^{5H-b}F'\big)\in\mathcal L .
\]
P85C4 now gives a finite integral cocycle representing
$p^{c_1}(y'+\Phi m)$. This is the third exponent in Q84T4.
% END P85C5

% BEGIN P85C6
\paragraph{P85C6. The ambiguity is an integral coboundary.}
Consider two potentials $j(\sigma)-b_0$, $j(\sigma')-b_0'$ of
the same $y$ on the very same module. Their Witt solutions satisfy
$m-m'\in T$ by P85C1; write this vector as $x$. Thus
$b_0-b_0'=1\otimes x$ by uniqueness of the crystalline value of
$\Phi(m-m')$. Their horizontal correction difference has
$\nabla(\sigma-\sigma')=0$, hence equals $\Xi(d)$ for its
constant term. This follows directly in the horizontal frame
because $u\,d/du$ kills only constants. Both corrections are
$\Delta$-invariant, so this period is invariant under the whole
group by the two factors of P85C4. Subtract the potential
coboundaries: the two resulting $V$-valued cocycles differ by
$g\mapsto(g-1)(-x)$ with $x$ in the integral lattice.

If $m_s\in M^\Delta$, the special vector $y_s=\Phi m_s-m_s$
has zero potential. Choose Witt solution $m_s$, correction
$\sigma_s=\Phi m_s$ and crystalline value $b_{0,s}=j(\Phi m_s)$.
The equation $\nabla\Phi m_s=\zeta(\Phi m_s-m_s)$ follows either
by the $p\Phi$ recurrence or directly by summing its telescoping
series. Thus adding $(\Phi-1)m_s$ to $y$ leaves its associated
cocycle unchanged. Addition and multiplication by $p^j$ commute
with all four defining identities: a common denominator for two
crystalline values is the product of their individual denominators.
Finally, a realized morphism sends the Witt equation to the Witt
equation of its image and takes $\sigma$ to its descent-invariant
image with derivative the corresponding $\zeta$. Both comparisons
on common coordinates commute with the lattice morphism, so
its tensor with $B_{\rm cris}$ sends the primitive to one for the
image. Taking Galois differences gives the actual postcomposed
cocycle. These facts prove functoriality and all the equivalences
used when comparing the two representatives in Q84T4.
% END P85C6
